% Options for packages loaded elsewhere
% Options for packages loaded elsewhere
\PassOptionsToPackage{unicode}{hyperref}
\PassOptionsToPackage{hyphens}{url}
\PassOptionsToPackage{dvipsnames,svgnames,x11names}{xcolor}
%
\documentclass[
  number,
  preprint]{elsarticle}
\usepackage{xcolor}
\usepackage{amsmath,amssymb}
\setcounter{secnumdepth}{5}
\usepackage{iftex}
\ifPDFTeX
  \usepackage[T1]{fontenc}
  \usepackage[utf8]{inputenc}
  \usepackage{textcomp} % provide euro and other symbols
\else % if luatex or xetex
  \usepackage{unicode-math} % this also loads fontspec
  \defaultfontfeatures{Scale=MatchLowercase}
  \defaultfontfeatures[\rmfamily]{Ligatures=TeX,Scale=1}
\fi
\usepackage{lmodern}
\ifPDFTeX\else
  % xetex/luatex font selection
\fi
% Use upquote if available, for straight quotes in verbatim environments
\IfFileExists{upquote.sty}{\usepackage{upquote}}{}
\IfFileExists{microtype.sty}{% use microtype if available
  \usepackage[]{microtype}
  \UseMicrotypeSet[protrusion]{basicmath} % disable protrusion for tt fonts
}{}
\makeatletter
\@ifundefined{KOMAClassName}{% if non-KOMA class
  \IfFileExists{parskip.sty}{%
    \usepackage{parskip}
  }{% else
    \setlength{\parindent}{0pt}
    \setlength{\parskip}{6pt plus 2pt minus 1pt}}
}{% if KOMA class
  \KOMAoptions{parskip=half}}
\makeatother
% Make \paragraph and \subparagraph free-standing
\makeatletter
\ifx\paragraph\undefined\else
  \let\oldparagraph\paragraph
  \renewcommand{\paragraph}{
    \@ifstar
      \xxxParagraphStar
      \xxxParagraphNoStar
  }
  \newcommand{\xxxParagraphStar}[1]{\oldparagraph*{#1}\mbox{}}
  \newcommand{\xxxParagraphNoStar}[1]{\oldparagraph{#1}\mbox{}}
\fi
\ifx\subparagraph\undefined\else
  \let\oldsubparagraph\subparagraph
  \renewcommand{\subparagraph}{
    \@ifstar
      \xxxSubParagraphStar
      \xxxSubParagraphNoStar
  }
  \newcommand{\xxxSubParagraphStar}[1]{\oldsubparagraph*{#1}\mbox{}}
  \newcommand{\xxxSubParagraphNoStar}[1]{\oldsubparagraph{#1}\mbox{}}
\fi
\makeatother


\usepackage{longtable,booktabs,array}
\usepackage{calc} % for calculating minipage widths
% Correct order of tables after \paragraph or \subparagraph
\usepackage{etoolbox}
\makeatletter
\patchcmd\longtable{\par}{\if@noskipsec\mbox{}\fi\par}{}{}
\makeatother
% Allow footnotes in longtable head/foot
\IfFileExists{footnotehyper.sty}{\usepackage{footnotehyper}}{\usepackage{footnote}}
\makesavenoteenv{longtable}
\usepackage{graphicx}
\makeatletter
\newsavebox\pandoc@box
\newcommand*\pandocbounded[1]{% scales image to fit in text height/width
  \sbox\pandoc@box{#1}%
  \Gscale@div\@tempa{\textheight}{\dimexpr\ht\pandoc@box+\dp\pandoc@box\relax}%
  \Gscale@div\@tempb{\linewidth}{\wd\pandoc@box}%
  \ifdim\@tempb\p@<\@tempa\p@\let\@tempa\@tempb\fi% select the smaller of both
  \ifdim\@tempa\p@<\p@\scalebox{\@tempa}{\usebox\pandoc@box}%
  \else\usebox{\pandoc@box}%
  \fi%
}
% Set default figure placement to htbp
\def\fps@figure{htbp}
\makeatother





\setlength{\emergencystretch}{3em} % prevent overfull lines

\providecommand{\tightlist}{%
  \setlength{\itemsep}{0pt}\setlength{\parskip}{0pt}}



 
\usepackage[]{natbib}
\bibliographystyle{elsarticle-num}


\makeatletter
\@ifpackageloaded{caption}{}{\usepackage{caption}}
\AtBeginDocument{%
\ifdefined\contentsname
  \renewcommand*\contentsname{Table of contents}
\else
  \newcommand\contentsname{Table of contents}
\fi
\ifdefined\listfigurename
  \renewcommand*\listfigurename{List of Figures}
\else
  \newcommand\listfigurename{List of Figures}
\fi
\ifdefined\listtablename
  \renewcommand*\listtablename{List of Tables}
\else
  \newcommand\listtablename{List of Tables}
\fi
\ifdefined\figurename
  \renewcommand*\figurename{Figure}
\else
  \newcommand\figurename{Figure}
\fi
\ifdefined\tablename
  \renewcommand*\tablename{Table}
\else
  \newcommand\tablename{Table}
\fi
}
\@ifpackageloaded{float}{}{\usepackage{float}}
\floatstyle{ruled}
\@ifundefined{c@chapter}{\newfloat{codelisting}{h}{lop}}{\newfloat{codelisting}{h}{lop}[chapter]}
\floatname{codelisting}{Listing}
\newcommand*\listoflistings{\listof{codelisting}{List of Listings}}
\makeatother
\makeatletter
\makeatother
\makeatletter
\@ifpackageloaded{caption}{}{\usepackage{caption}}
\@ifpackageloaded{subcaption}{}{\usepackage{subcaption}}
\makeatother
\journal{Reliability Engineering & System Safety}
\usepackage{bookmark}
\IfFileExists{xurl.sty}{\usepackage{xurl}}{} % add URL line breaks if available
\urlstyle{same}
\hypersetup{
  pdftitle={A Benchmark-Integrity Audit Protocol for Conformal Battery State-of-Health Evaluation: Recent-Label Dependence, Schema-Bundle Sensitivity, and Threshold-Decision Risk},
  pdfauthor={Qingsong Shan; Qianning Liu},
  pdfkeywords={conformal prediction, uncertainty quantification, battery
state-of-health, prognostics and health management, predictive
maintenance, benchmark integrity, distribution shift},
  colorlinks=true,
  linkcolor={blue},
  filecolor={Maroon},
  citecolor={Blue},
  urlcolor={Blue},
  pdfcreator={LaTeX via pandoc}}


\setlength{\parindent}{6pt}
\begin{document}

\begin{frontmatter}
\title{A Benchmark-Integrity Audit Protocol for Conformal Battery
State-of-Health Evaluation: Recent-Label Dependence, Schema-Bundle
Sensitivity, and Threshold-Decision Risk}
\author[1]{Qingsong Shan%
%
\fnref{fn1}}
 \ead{shanqingsong@jxufe.edu.cn} 
\author[1]{Qianning Liu%
\corref{cor1}%
\fnref{fn2}}
 \ead{liuqianning@jxufe.edu.cn} 

\affiliation[1]{organization={Jiangxi University of Finance and
Economics, School of Statistics and Data Science, Key Laboratory of Data
Science in Finance and
Economics},city={Nanchang},country={China},countrysep={,},postcodesep={}}

\cortext[cor1]{Corresponding author}
\fntext[fn1]{Open Researcher and Contributor Identifier (ORCID):
0000-0001-9231-5699}
\fntext[fn2]{Open Researcher and Contributor Identifier (ORCID):
0000-0002-7654-6451}
        
\begin{abstract}
Conformal intervals for lithium-ion battery state-of-health (SOH) inform
maintenance decisions, but cross-dataset validation can conceal
measurement conditions. We develop a six-gate benchmark-integrity audit
of prediction-time information, split integrity, cell-block uncertainty,
schema controls, threshold decisions, and localization sensitivity.
Twelve analyses use 13,831 eligible records from 13,876 CALCE, NASA, and
Oxford observations. A pooled-domain, cell-held-out baseline reaches
94.8\% coverage (cell-block lower bound 94.2\%). Across six
cross-dataset tasks, method-matched removal of the most recent measured
SOH yields statistically supported coverage reductions in 5 tasks, by as
much as 60.4 percentage points; the remaining task and pooled-domain
baseline are compatible with no effect. The as-used schema bundle
changes leave-NASA coverage by -12.0 percentage points and false
acceptance by +100.0 percentage points at the illustrative SOH cutoff
0.90 (exact paired-cell p=9.54e-07); an independent control bounds this
as a bundle-level association. A localized-residual comparator recovers
one point-coverage failure but establishes no shift-robust guarantee.
Each gate couples evidence requirements to claim bounds, yielding an
executable pre-submission check demonstrated on three public battery
benchmarks; detection performance remains unestablished.
\end{abstract}





\begin{keyword}
    conformal prediction \sep uncertainty quantification \sep battery
state-of-health \sep prognostics and health management \sep predictive
maintenance \sep benchmark integrity \sep 
    distribution shift
\end{keyword}
\end{frontmatter}
    

\section{Introduction}\label{introduction}

Data-driven prognostics and health management increasingly supplies the
evidence on which maintenance and serviceability decisions rest, and
closing the gap between research practice and trustworthy deployment is
a recognized priority for the reliability community
\citep{zio2022prognostics, lin2026recent}. Lithium-ion battery
state-of-health (SOH) estimation is a flagship case: high-impact
data-driven estimators are routinely evaluated on public cycle-level
datasets \citep{wang2023improved, meng2023long}, uncertainty
quantification is now treated as a first-class evaluation criterion for
such models \citep{basora2025benchmark, nemani2023uncertainty}, and
conformal prediction --- with its finite-sample coverage guarantee ---
is entering reliability-engineering and battery pipelines directly
\citep{xu2026uncertaintyaware, javanmardi2023conformal, soon2026enhancing}.
Because SOH intervals feed threshold decisions about maintenance,
replacement, and second-life use, interval miscalibration has direct
decision consequences
\citep{mitici2023dynamic, zhuang2023prognostic, wang2026multiaction}.

SOH uncertainty estimates are often evaluated through cross-dataset
benchmarks, where a model calibrated on one collection is tested on
another --- a realistic proxy for the domain shift that deployed
prognostics faces \citep{nejjar2024domain}. In maintenance- and
reliability-facing uses, it is tempting to read those intervals as
evidence that threshold decisions remain trustworthy under protocol
shift. Yet dataset composition and feature choices, rather than model
quality, can drive reported battery-prognostics performance
\citep{nagulapati2021capacity, geslin2023selecting}, and the reliability
of machine-learning evidence is itself an engineering quantity that
requires pre-deployment assessment \citep{chen2024evaluating}. This
paper argues that the cross-dataset reading is not warranted unless the
benchmark itself is audited as a measurement instrument and the
information available at prediction time is stated explicitly.

This paper consolidates that audit into an explicit, governed protocol
and demonstrates it on real cycle-level SOH records from the Center for
Advanced Life Cycle Engineering (CALCE), the National Aeronautics and
Space Administration (NASA), and Oxford. The demonstration bounds
external validity to these collections and centers a
reliability-engineering question that current evaluation practice leaves
implicit: how much of the apparent validity of conformal SOH evidence
depends on access to the most recent measured SOH, and what remains
interpretable after schema, protocol, and decision-level checks?

Two bodies of work bear directly on this question, but neither audits
the benchmark itself. Conformal methods in reliability engineering and
battery state estimation concentrate on producing or adapting predictive
intervals
\citep{xu2026uncertaintyaware, wang2025robust, javanmardi2023conformal, piao2025crulp, soon2026enhancing},
while transfer and domain-generalization methods concentrate on making
estimators survive distribution shift
\citep{nejjar2024domain, nejjar2026uncertaintyguided, fei2026decoupled}.
Non-exchangeability theory bounds how far conformal coverage can drift
\citep{barber2023beyond, oliveira2024split}, but it does not identify
which measurement-validity channel a particular benchmark contains.
Within the literature reviewed in Section \ref{sec-related-work}, no
integrated protocol treats a conformal SOH benchmark itself as the
object of reliability analysis. This is the gap addressed here.

The main contributions of this paper are as follows:

\begin{enumerate}
\def\labelenumi{\arabic{enumi}.}
\tightlist
\item
  It defines a benchmark-integrity audit protocol whose six governed
  gates cover the prediction-time information set, split integrity,
  cell-aware uncertainty, schema negative controls, threshold-decision
  consequences, and localization sensitivity.
\item
  It quantifies recent-SOH access as a load-bearing, task-dependent
  measurement condition, separating a method-selection ablation from
  method-matched, row-matched contrasts.
\item
  It audits as-used schema-bundle sensitivity with multiplicity-adjusted
  intervals and an independently permuted-feature negative control,
  showing how the control gate confines bundle effects to a bounded
  sensitivity warning instead of a spurious field-specific mechanism
  claim.
\item
  It reports a corrected leave-NASA (LNAS) schema-bundle
  false-acceptance increase of +100.0 pp at threshold 0.90 (exact
  paired-cell p = 9.54e-07; 21 effective cells and 1127 records),
  linking interval behavior to threshold-decision risk under the
  bundled, noncausal interpretation.
\item
  It bounds residual sensitivity with a localized-residual comparator
  envelope, distinguishing empirical recovery from a formal
  covariate-shift guarantee.
\end{enumerate}

The intended takeaway is operational: raw cross-dataset conformal
coverage becomes battery SOH reliability evidence only after the
benchmark's prediction-time information, split construction, controls,
and decision consequences have been audited --- and the protocol
developed here makes that audit concrete and executable.

The remainder of the paper is organized as follows.
Section~\ref{sec-related-work} reviews related work on conformal
prediction, battery SOH uncertainty, and benchmark integrity.
Section~\ref{sec-methodology} defines the audit design, its provenance,
and its claim boundaries. Section~\ref{sec-results} reports evidence on
recent-label dependence, schema-bundle sensitivity, residual drift,
decision consequences, and localized calibration.
Section~\ref{sec-discussion} discusses implications and limitations, and
Section~\ref{sec-conclusion} gives concluding remarks.

\section{Related Work}\label{sec-related-work}

\subsection{Conformal prediction: guarantees and their boundary
conditions}\label{conformal-prediction-guarantees-and-their-boundary-conditions}

Conformal prediction provides finite-sample predictive-coverage
guarantees under exchangeability
\citep{vovk2005algorithmic, lei2018distributionfree, angelopoulos2023conformal}.
Conformalized quantile regression calibrates quantile-regression
estimates to retain finite-sample coverage
\citep{romano2019conformalized}, and weighted conformal variants extend
the guarantee to covariate shift when the shift mechanism can be
estimated \citep{tibshirani2019covariate}. The boundary conditions of
these guarantees are now well mapped. When calibration and test data are
dependent or non-exchangeable, coverage deviates by an amount that
theory can bound but not eliminate
\citep{barber2023beyond, oliveira2024split}, and ensemble batch
prediction intervals extend conformal thinking to dependent time series
\citep{xu2023conformal}. Online and adaptive variants track arbitrary
distribution shifts, but they require sequential feedback after each
prediction \citep{gibbs2021adaptive, gibbs2024conformal}, which a
one-shot cross-dataset evaluation does not provide. A parallel line
sharpens the distinction between marginal and conditional validity:
exact conditional coverage is unattainable in general, and localized or
reweighted calibration offers only partial remedies
\citep{gibbs2025conformal, guan2022localized, hore2024conformal}.
Conformal risk control generalizes coverage to other monotone risks,
including decision-relevant error rates, with non-exchangeable
extensions
\citep{angelopoulos2024conformalrisk, farinhas2024nonexchangeable}.

These guarantees do not by themselves certify a reliability benchmark.
If the information set includes recent labels, dataset identity, or
bookkeeping fields unavailable in the intended decision setting, nominal
coverage can describe the constructed benchmark rather than the intended
threshold-decision problem. The localized-residual analysis used in this
paper is a diagnostic inspired by this literature, not an implementation
of density-ratio weighted conformal prediction.

\subsection{Uncertainty-aware battery SOH and remaining useful life
prognostics}\label{uncertainty-aware-battery-soh-and-remaining-useful-life-prognostics}

Battery SOH and remaining useful life (RUL) studies have long reported
probabilistic forecasts and prognostic performance under uncertainty
\citep{saxena2010evaluating, he2011remaining, liu2020futurecapacity},
and real-application SOH estimation has a mature methodological
literature \citep{berecibar2016critical}. Within the
reliability-engineering literature specifically, data-driven battery
prognostics is a flagship application: high-impact estimators define the
benchmark conventions on CALCE-, NASA-, and Oxford-style cycle datasets
\citep{wang2023improved, meng2023long, wei2023prediction}, and an
explicitly uncertainty-aware thread now reports quantile intervals
validated on the same public datasets \citep{zhang2024reliability},
Bayesian-calibrated SOH uncertainty under deployment shift
\citep{zhu2025bayesian}, probabilistic degradation-trajectory health
assessment \citep{che2024predictive, lin2023battery}, and recent
dynamical-model variants \citep{wang2025remaining}. Label scarcity at
deployment is a recognized constraint in this literature: adaptive
prediction with few or no in-service capacity labels is studied as its
own problem \citep{zhang2022adaptive, pan2024probabilistic}, which makes
the availability of a recent measured SOH a substantive modeling
assumption rather than a technicality. Dedicated benchmarks and reviews
now treat uncertainty quantification as a first-class evaluation
criterion for deep-learning prognostics
\citep{basora2025benchmark, lin2026recent, nemani2023uncertainty, salinascamus2025comprehensive}.

Conformal prediction itself has entered this application class in three
roles. It produces calibrated RUL and degradation intervals
\citep{xu2026uncertaintyaware, liu2026credible, javanmardi2023conformal, piao2025crulp, soon2026enhancing},
adapts interval inference to corrupted inputs or engineering-design
uncertainty \citep{wang2025robust, wu2026sequential}, and supports
anomaly detection or width-aware calibration
\citep{yuan2026conformal, xu2026widthoptimized}. Related industrial
time-series work also documents the exchangeability violations that
motivate such adaptations \citep{zhang2025uncertainty}. These methods
improve or diagnose interval estimators, but they largely take the
evaluation benchmark as given. That presupposition is precisely what the
present audit tests.

\subsection{Cross-dataset evaluation, transfer, and domain shift in
prognostics and health
management}\label{cross-dataset-evaluation-transfer-and-domain-shift-in-prognostics-and-health-management}

Transfer research likewise pursues several remedies for distribution
shift. Domain-adaptation methods move SOH or RUL estimators across
batteries, operating profiles, and laboratory or field domains
\citep{ye2021domainadversarial, vilsen2022field, shen2024review, nejjar2024domain, nejjar2026uncertaintyguided, wang2026twostage}.
Condition-aware and domain-generalization methods target fleet or
protocol variation \citep{zhang2023multitask, fei2026decoupled}, while
adjacent reliability work mitigates shift in data-driven risk assessment
\citep{wang2025mitigating}. This literature is remedy-side: it seeks
models that survive shift and treats the benchmark schema as a fixed
prerequisite for comparison. The present paper takes the complementary
audit-side view: before asking whether a model transfers, it asks
whether the benchmark's own construction --- its information set,
splits, feature schema, and decision endpoints --- supports the
reliability conclusion being drawn.

\subsection{Benchmark integrity, leakage, and evaluation
validity}\label{benchmark-integrity-leakage-and-evaluation-validity}

Machine-learning reliability research warns that apparent predictive
performance can be driven by unintended signals, undocumented data
provenance, or data-splitting and feature-availability choices
\citep{gebru2021datasheets, geirhos2020shortcut, kapoor2023leakage}. The
battery and PHM literatures now contain concrete instances: feature
selection can produce misleading battery-lifetime predictions
\citep{geslin2023selecting}, training-set size and composition ---
rather than model class --- can drive reported capacity-estimation
accuracy \citep{nagulapati2021capacity}, explainability analysis reveals
which features an SOH model actually exploits
\citep{wang2023explainabilitydriven}, and open benchmark suites codify
the feature schemas whose integrity this paper audits
\citep{wang2024open}. Within reliability engineering, an
evaluation-validity genre is emerging: simple baselines can expose what
domain-shift benchmarks actually measure in condition monitoring
\citep{panic2027simple}, and pre-deployment reliability assessment of
machine-learning evidence is treated as its own engineering problem
\citep{chen2024evaluating}.

This paper connects these threads. Within the literature reviewed above,
we did not find an integrated audit protocol that treats a conformal SOH
benchmark itself as the object of reliability analysis. The unit of
analysis here is the benchmark evidence --- how conformal interval
behavior and threshold decisions change across recent-label, split,
schema-bundle, residual-drift, and SOH-threshold diagnostics. The
protocol integrates these diagnostics into six gates, each with an
explicit evidentiary requirement and a consequence that bounds the
admissible claim when the gate fails.

\section{Methodology}\label{sec-methodology}

The methodological contribution of this paper is a governed
benchmark-integrity audit protocol for conformal SOH evidence. Table
\ref{tbl-audit-protocol} defines its six gates: each names the
measurement condition examined, the evidence required, and the
consequence of failure. The protocol consolidates a sequence of
completed diagnostic analyses whose component decision rules were
documented in their corresponding experiment plans before those analyses
were executed. The six-gate synthesis was developed after that
diagnostic program and is not presented as a prospectively registered
protocol. A gate failure bounds the resulting claim rather than
terminating the analysis; this design allows negative and heterogeneous
findings to remain informative instead of being discarded. The remainder
of this section instantiates the protocol on real cross-dataset battery
evidence.

\begin{table}

\caption{\label{tbl-audit-protocol}Governed benchmark-integrity audit protocol for conformal SOH evidence. Each gate names the measurement condition examined, the evidence required, and the consequence of failure applied in this audit.}

\centering{

\centering
\scriptsize


\begin{tabular}{>{\raggedright\arraybackslash}p{0.22\linewidth}>{\raggedright\arraybackslash}p{0.46\linewidth}>{\raggedright\arraybackslash}p{0.24\linewidth}}
\toprule
Gate & Evidence required & Consequence of failure \\
\midrule
G1 Information set & State whether recent measured SOH or another oracle label is available at prediction time & No claim for settings without recent labels \\
G2 Split integrity & Verify task-appropriate disjoint training, validation, calibration, and test indices; require cell-disjoint histories where the task claims cell-level holdout & Treat coverage as split-contingent \\
G3 Uncertainty units & Report cell-block uncertainty rather than row-level iid intervals & Treat uncertainty as optimistic \\
G4 Schema sensitivity & Compare schema contrasts with an independent permuted-feature control and multiplicity adjustment & Bundle-level warning; no mechanism claim \\
G5 Decision consequences & Report threshold false acceptance and false rejection with record denominators, effective cell counts, and sparse-cell flags & Do not promote coverage to serviceability evidence \\
G6 Localization sensitivity & Compare standard and localized calibration as a sensitivity envelope & Failure: no empirical recovery; apparent success: point-coverage sensitivity only. Neither outcome authorizes a shift-robust claim without a guarantee-bearing method and verified assumptions \\
\bottomrule
\end{tabular}

}

\end{table}%

The reported evidence is organized around governed diagnostic analyses
rather than selective numerical filtering. A scripted results table
aggregates the outputs used for the manuscript, and the reported values
are taken from that table. All active rows are real battery
state-of-health (SOH) records from CALCE, NASA, and Oxford
\citep{calce_battery_data, nasa_pcoe_battery_data, howey2017oxford}; no
synthetic records are included in the reported evidence. The audited
object is the benchmark evidence produced by these public cycle-level
datasets; the protocol itself, not a new battery-management algorithm,
is the methodological contribution. The 12 completed diagnostic analyses
are a pooled-domain, cell-held-out coverage baseline, a recent-label
ablation, a cross-protocol transfer diagnostic, a residual-bias screen
(tests whether calibration residual signs vary systematically with cycle
position), a host-timing screen (tests whether results shift with the
wall-clock execution host), a negative-control check, a method-family
comparison, a threshold decision-utility audit, a schema-cleaning audit,
a hard-regime audit, the schema-to-decision audit, and a
localized-residual comparator. Supplementary Tables S1--S2 map every
analysis to its gate, task, feature schema, method family, parameter
grid, selection rule, seed, plan identifier, governed rerun, and
deviation record. The archived plans and regeneration sources are fixed
at public replication commit
\texttt{42f1120ed1f5cbc0282a706ebd4bec24a8d19d6e}.

The physical case spans three laboratory platforms rather than one
nominally homogeneous fleet. The CALCE subset contains six CS2 and six
CX2 small prismatic cells with LiCoO\(_2\) cathodes and capacity ratings
of 1.10 and 1.35 Ah. The archive documents
constant-current/constant-voltage charging for CS2 and 0.5C or 1C
constant-current discharge for the selected CS2/CX2 cells; the retained
row files do not encode ambient temperature \citep{calce_battery_data}.
The NASA subset contains 25 Li-ion cells from multiple Prognostics
Center of Excellence batches; the repository describes charge,
discharge, and impedance profiles for nominal 2-Ah cells, and the
retained metadata spans 4, 22, 24, and 43 \(^\circ\)C, but does not
consistently encode cathode chemistry \citep{nasa_pcoe_battery_data}.
Oxford contributes eight 740-mAh Kokam pouch cells aged in a 40
\(^\circ\)C chamber with constant-current/constant-voltage charging and
an ARTEMIS urban-drive discharge profile \citep{howey2017oxford}. These
differences make label and protocol comparability part of the audit
boundary rather than a solved preprocessing detail. Table
\ref{tbl-datasets} summarizes the 45-cell source evidence base of 13,876
cycle rows and the 13,831 rows eligible for a strictly prior recent-SOH
feature.

\begin{table}

\caption{\label{tbl-datasets}Real battery datasets after preprocessing. Rows are source/eligible; every eligible row has an earlier same-cell SOH record.}

\centering{

\centering
\scriptsize


\begin{tabular}{l>{\raggedright\arraybackslash}p{0.27\linewidth}>{\raggedright\arraybackslash}p{0.31\linewidth}rcc}
\toprule
Dataset & Cell platform & Cycling / temperature & Cells & Rows & SOH range \\
\midrule
CALCE & CS2/CX2 prismatic, LiCoO$_2$, 1.10/1.35 Ah & CS2: CC--CV charge; selected CS2/CX2: 0.5C/1C discharge; temperature absent & 12 & 11{,}604/11,592 & 0.40--1.01 \\
NASA & Multi-batch Li-ion cells; archive nominal 2 Ah; cathode not consistently encoded & Charge/discharge/EIS; retained metadata span 4--43 $^\circ$C & 25 & 1{,}753/1,728 & 0.41--1.00 \\
Oxford & Kokam SLPB533459H4 pouch, 0.74 Ah & CC--CV charge; ARTEMIS urban-drive discharge; chamber 40 $^\circ$C & 8 & 519/511 & 0.62--1.00 \\
\midrule
Total & Three public laboratory platforms & Heterogeneous protocols and measurement conventions & 45 & 13,876/13,831 & 0.40--1.01 \\
\bottomrule
\end{tabular}

}

\end{table}%

The seven evaluation tasks use the same model families but change the
target-domain role. Leave-domain-out tasks fit, select, and calibrate on
the other two domains and test on all eligible rows of the named
held-out domain. Target-recalibrated tasks add early target-domain rows:
progress 0--20\% enters training, 20--35\% calibration, 35--55\%
validation, and rows after 55\% form the test set. These boundaries are
pairwise disjoint. Table \ref{tbl-task-design} reports the resulting row
counts; the held-out target contains 12 CALCE, 25 NASA, or 8 Oxford
cells, while Main uses 10 test cells. Main is a pooled-domain
within-benchmark baseline: all three datasets contribute to fitting and
test sets, while cell identities remain disjoint; it is not a
single-physical-protocol validation.

\begin{table}

\caption{\label{tbl-task-design}Operational task construction and row counts for training, calibration, validation, and testing.}

\centering{

\centering
\scriptsize


\begin{tabular}{llrl}
\toprule
Task & Construction & Test cells & Train/Cal/Val/Test rows \\
\midrule
Main & Cell-disjoint mixed-domain held-out split & 10 & 7961/880/680/2323 \\
LCAL & Other domains to all CALCE rows & 12 & 1122/241/116/11592 \\
TCAL & Other domains + early CALCE to late CALCE & 12 & 4566/1754/2344/5167 \\
LNAS & Other domains to all NASA rows & 25 & 7082/664/636/1728 \\
TNAS & Other domains + early NASA to late NASA & 25 & 12447/279/322/783 \\
LOxf & Other domains to all Oxford rows & 8 & 7718/855/608/511 \\
TOxf & Other domains + early Oxford to late Oxford & 8 & 13421/71/106/233 \\
\bottomrule
\end{tabular}

}

\end{table}%

The Main row counts follow governed cell partitions and the progress
windows specified below. Leave-domain tasks reuse those partitions after
restricting training, calibration, and validation to non-target domains;
consequently, their source-domain rows need not sum to every eligible
source-domain row in Table \ref{tbl-datasets}. Target-recalibrated
boundaries are time-forward and pairwise row-disjoint.

Main is a seed-42, domain-stratified cell split, not a within-cell
forecasting split. Within each domain, shuffled cells are assigned
55\%/20\%/10\%/remainder to training/calibration/validation/test;
calibration retains progress at most 0.55 and test retains progress at
least 0.35. The four cell sets are disjoint, and no test cell
contributes an earlier row to model fitting, selection, or calibration.
The complete cell membership is recorded below; the governed split
manifest preserves row indices and its checksum.

\begin{quote}
\scriptsize
\textbf{Complete Main split membership} (domain entries are row counts).\par
\textbf{Train} (7,961 rows; CALCE 6,839; NASA 879; Oxford 243): CALCE\_CS2\_33, CALCE\_CS2\_36, CALCE\_CX2\_33, CALCE\_CX2\_34, CALCE\_CX2\_36, CALCE\_CX2\_38, NASA\_B0007, NASA\_B0025, NASA\_B0026, NASA\_B0028, NASA\_B0029, NASA\_B0031, NASA\_B0042, NASA\_B0043, NASA\_B0048, NASA\_B0053, NASA\_B0054, NASA\_B0055, NASA\_B0056, Oxford\_Cell1, Oxford\_Cell3, Oxford\_Cell4, Oxford\_Cell6.

\textbf{Cal} (880 rows; CALCE 639; NASA 216; Oxford 25): CALCE\_CS2\_35, CALCE\_CS2\_38, NASA\_B0006, NASA\_B0027, NASA\_B0045, NASA\_B0046, NASA\_B0047, Oxford\_Cell5.

\textbf{Val} (680 rows; CALCE 564; NASA 44; Oxford 72): CALCE\_CS2\_37, NASA\_B0032, NASA\_B0051, Oxford\_Cell2.

\textbf{Test} (2,323 rows; CALCE 1,968; NASA 255; Oxford 100): CALCE\_CS2\_34, CALCE\_CX2\_35, CALCE\_CX2\_37, NASA\_B0005, NASA\_B0018, NASA\_B0030, NASA\_B0044, NASA\_B0049, Oxford\_Cell7, Oxford\_Cell8.
\end{quote}

For cell \(i\) and recorded cycle \(t\), \texttt{prev\_soh} is
\(SOH_{i,t-1}\): the immediately preceding recorded same-cell
capacity-SOH after sorting by dataset, cell identifier, and cycle index.
The \(t-1\) measurement must occur before the prediction event at \(t\).
Each cell's first row has no such measurement and is excluded before
every split is constructed (45 rows in total); no same-row target fill
or other imputation is used. Thus recent-SOH is an oracle only in the
sense that it assumes a previous capacity measurement is actually
available at prediction time, not because it accesses \(SOH_{i,t}\).

The primary nominal coverage level is 90\%. The manuscript reports
coverage, width, dependence-aware uncertainty, threshold decision
utilities, schema diagnostics, and localized-residual comparator
outcomes from the scripted results table. Before the schema-to-decision
analysis was executed, its experiment plan defined five contrasts across
seven tasks and seven endpoints: one coverage endpoint plus the
predeclared SOH thresholds 0.65, 0.70, 0.75, 0.80, 0.85, 0.90, giving
245 tests. Coverage and false-acceptance deltas in Tables
\ref{tbl-schema-deltas} and \ref{tbl-false-acceptance-deltas} use
Bonferroni-adjusted, SE-based normal intervals from 1000 cell-block
resamples; the family-wise level 0.05 gives a per-test level near
\(2\times10^{-4}\) and z \(\approx\) 3.7. Zero-SE boundaries use the
exact paired-cell sign test. By contrast, the Main 94.2\%--97.1\%
interval and the lower bounds in Table \ref{tbl-adaptive-envelope} are
2.5/97.5-percentile cell-block bootstrap intervals, while Table
\ref{tbl-decision-tradeoff} reports uncorrected 95\% cell-block normal
intervals. The limited number of battery-cell blocks makes extreme-tail
normal widths uncertain; this is an additional reason to treat coverage
flags as sensitivity signals and require a negative-control gate. The
same data interface generates text, tables, and captions to avoid manual
transcription.

The nominal target is used as the evaluation criterion, while individual
conformal variants select their calibration quantiles strictly on a
disjoint validation split (never tuned on the test set). For the
pooled-domain, cell-held-out diagnostic baseline, the selected operating
point uses q-level 0.93. This yields a realized test coverage of 94.8\%;
the benign validation-to-test over-coverage is not treated as a new
performance claim, but as the baseline's `apparent reliability' before
domain shift or artifact removal, setting the reference for subsequent
measurement-validity checks relative to the 90\% coverage target.

Persistence-anchor behavior is used as a diagnostic probe because it
exposes recent-label and anchor dependence. It is not presented as a
deployable or superior method. Cross-protocol failures are interpreted
as validity evidence, not as a failed attempt to force a positive
transfer claim. The analysis therefore treats negative and heterogeneous
findings as informative: a failure after schema cleaning is evidence
consistent with residual protocol drift, while a small, directionally
heterogeneous, or control-inconsistent schema effect bounds the claim
rather than generalizing it. Residual drift is retained as a boundary
that is not fully separable from recent-label dependence in this
evidence base.

The audit examines four diagnostic dimensions:

\begin{enumerate}
\def\labelenumi{\arabic{enumi}.}
\tightlist
\item
  recent-label / anchor dependence;
\item
  as-used schema-bundle sensitivity;
\item
  residual protocol drift after schema checks;
\item
  SOH serviceability-threshold decision risk under finite-cell evidence.
\end{enumerate}

The Main pooled-domain, cell-held-out diagnostic baseline is the
selected persistence-anchor protocol Mondrian conformal predictor
\citep{vovk2005algorithmic}; the split-conformal constructions and
coverage definitions follow standard practice
\citep{angelopoulos2023conformal}. The diagnostic families include
quantile regression (QR), conformalized quantile regression (CQR)
\citep{romano2019conformalized}, and natural-gradient boosting (NGBoost)
\citep{duan2020ngboost}. For the no-recent-SOH ablation condition, the
selected diagnostic is a gradient-boosting domain-Mondrian conformal
variant, a group-conditional calibration construction suited to settings
where exchangeability across domains is not credible
\citep{barber2023beyond}; it is used only to quantify the dependence of
the benchmark on recent-label access and is not presented as a method
contribution. Because the calibration-selection rule, not the analyst,
chooses the per-condition operating point, the full-feature and
no-recent-SOH conditions can select different conformal families; the
resulting pooled-domain coverage change therefore reflects the selection
rule's response to removing recent labels rather than a same-method
feature ablation, and the method-matched feature-isolation contrast is
supplied separately by the row-matched recent-label reference comparison
reported in Section \ref{sec-results}. Gradient-boosting CQR is used for
decision-utility diagnostics, schema comparisons, and the
localized-residual comparator. Training, validation, calibration, and
test indices are pairwise disjoint per task. The as-used bundle contains
\texttt{source\_row\_id}, \texttt{target\_cycle\_index}, and the
permuted control column, whereas the clean-original schema removes all
three; the bundle contrast therefore cannot isolate a
bookkeeping-specific effect. The independent control adds a fixed-seed
(seed 1033) permutation of \texttt{log\_cycle\_index} to the clean
schema and refits the complete pipeline on the same rows. This control
uses one permutation draw (\(B=1\)), not a sampled permutation null
distribution; it is therefore used only as a veto on mechanism
attribution.

A decision-risk cell is defined relative to a capacity-SOH
serviceability threshold and a predicted SOH interval with lower bound L
and upper bound U for a test record with true SOH y. The record is
threshold-acceptable when its lower bound L is at or above the
threshold, threshold-unacceptable when its upper bound U is below the
threshold, and uncertain when the interval straddles the threshold. A
false-acceptance event is a threshold-acceptable record whose true SOH y
is below the threshold; a false-rejection event is a
threshold-unacceptable record whose true SOH y is at or above the
threshold. The reported false-acceptance rate divides false-acceptance
events by the number of actually threshold-unacceptable records, the
false-rejection rate divides false-rejection events by the number of
actually threshold-acceptable records, and the uncertain rate divides
uncertain records by all records; each rate carries a cell-block
bootstrap confidence interval that resamples battery cells with
replacement. When every contributing cell block contains zero events,
that bootstrap has zero empirical variance; the table therefore marks
the interval as bootstrap-degenerate rather than displaying a zero-width
confidence interval. Such a boundary is not evidence of zero underlying
risk. Auditing these monotone decision risks alongside coverage follows
the same motivation as conformal risk control, which generalizes
coverage guarantees to decision-relevant error rates
\citep{angelopoulos2024conformalrisk}. These thresholds are
capacity-based decision cutoffs, not demonstrated electrochemical or
thermal safety boundaries. To avoid over-reading sparse cells, a
small-denominator suppression rule flags any decision cell whose
denominator is below 5 as exploratory rather than as a headline rate.

Schema-bundle sensitivity is adjudicated by a rule documented before the
schema-to-decision analysis and applied identically to every task. A
task is flagged when the as-used-minus-clean coverage interval excludes
zero and the absolute delta is at least 3 pp.~This flag is necessary but
not sufficient for a mechanistic claim: the independent permuted-feature
control must remain null across both its coverage contrast and all six
false-acceptance endpoints. The bundle's false-acceptance-rate delta is
evaluated separately with multiplicity adjustment or, for
zero-standard-error boundaries, an exact paired-cell sign test; cells
below 5 effective observations are exploratory. This hierarchy prevents
a bundled association from being relabeled as a bookkeeping mechanism,
while still allowing a corrected bundle-level decision association to be
reported without field-specific attribution.

Two dependence boundaries are explicit. First, the conformal calibration
objects are cycle-level residuals from longitudinal cell records, so the
usual exchangeability interpretation is conditional on the task
construction rather than a free guarantee for independent operating
histories; under such dependence, split-conformal coverage is expected
to deviate from nominal by an amount that theory bounds but does not
eliminate \citep{oliveira2024split, xu2023conformal}. The uncertainty
intervals therefore resample cells as blocks and are read as
dependence-aware diagnostics, but the present evidence does not replace
a fully cell-independent prospective forecasting design. Second,
recent-SOH is treated as an oracle recent label in the main dependence
probe: it may be a legitimate autoregressive measurement in a laboratory
setting where the previous capacity has just been measured, but it is
not automatically available in an online battery-management setting. The
audit therefore interprets recent-SOH access as a measurement condition
to be declared, not as leakage in every possible deployment scenario.

The SOH label is capacity-normalized within the harmonized project data
interface, following the conventions of published SOH benchmark suites
on these datasets \citep{wang2024open, zhang2024reliability}, and Table
\ref{tbl-datasets} reports the ranges after quality assurance (QA).
Because the source datasets differ in cycling protocol, temperature
range, and capacity measurement practice, the analysis treats
cross-dataset label comparability as part of the benchmark-validity
boundary. The current manuscript does not establish a full
label-harmonization audit over alternative denominators, smoothing
rules, or anomaly-removal policies; such sensitivity analysis would be
needed before attributing every cross-dataset coverage failure to
physical protocol drift or model behavior.

Because a standard conformal failure is incomplete without a
localization sensitivity check, the audit also evaluates a heuristic
localized-residual CQR comparator in the spirit of localized conformal
prediction \citep{guan2022localized}. Calibration and target features
are standardized by calibration means and standard deviations. For
target row \(x\), calibration score \(r_i\) receives Gaussian weight
\(w_i(x)=\exp\{-d_i(x)^2/(2h^2)\}+10^{-12}\), where \(d_i^2\) is the
mean squared standardized distance over the recorded common shift
features. The local radius is the weighted conformal quantile. Bandwidth
\(h\in\{0.25,0.50,1,2,4,\infty\}\) is selected only on validation data
by the predeclared coverage-first, then width, objective; test metrics
are not used. This is motivated by covariate-shift conformal work
\citep{tibshirani2019covariate}, but it does not estimate a density
ratio and carries no formal covariate-shift guarantee. Meeting the
envelope means that point coverage reaches nominal; a lower-bound
shortfall crosses only when the standard percentile lower bound is below
nominal and the localized lower bound is at or above it. The result is
not interpreted as method superiority.

\section{Results}\label{sec-results}

This section reports what the audit protocol detects when applied to a
cross-dataset conformal SOH benchmark; the engineering object is the
reliability evidence produced by the benchmark, not a new SOH estimator.
The decision question is how interval coverage and SOH threshold
decisions change across recent-label access, split construction,
feature-schema sensitivity, residual protocol drift, and localized
residual weighting diagnostics.

The evidence base is a set of diagnostic contrasts over the same
real-data benchmark splits. Task abbreviations use Main for the
pooled-domain, cell-held-out task, leave-CALCE (LCAL), target-CALCE
(TCAL), leave-NASA (LNAS), target-NASA (TNAS), leave-Oxford (LOxf), and
target-Oxford (TOxf). Table \ref{tbl-evidence-map} reports the
argument-bearing contrasts needed to interpret reliability evidence;
ancillary timing and residual-stress screens are treated as robustness
checks rather than headline findings.

\begin{table}

\caption{\label{tbl-evidence-map}Gate-by-gate audit ledger. Later tables report interval estimates. Pass* denotes complete reporting with bounded interpretation; pp denotes percentage points.}

\centering{

\centering
\scriptsize


\begin{tabular}{>{\raggedright\arraybackslash}p{0.05\linewidth}>{\raggedright\arraybackslash}p{0.50\linewidth}>{\raggedright\arraybackslash}p{0.10\linewidth}>{\raggedright\arraybackslash}p{0.25\linewidth}}
\toprule
Gate & Case evidence & Verdict & Claim bound \\
\midrule
G1 & Prior SOH declared; removal -60.4 pp in TCAL; TOxf +4.3 pp with interval including zero & Conditional & Exclude settings without prior labels \\
G2 & Main cell-disjoint, governed source splits for leave-domain tasks, forward disjoint recalibration rows & Pass & Split-contingent \\
G3 & Cell-block intervals, effective task count 8--25 cells & Pass & No row-iid precision \\
G4 & LCAL/LNAS bundle effects, LNAS -12.0 pp, control non-null in LCAL/TCAL & Fail & Bundle only, no mechanism \\
G5 & Both threshold errors, uncertainty status, cell/record counts, no sparse displayed cell & Pass* & Illustrative, not deployment rates \\
G6 & TNAS recovers one point failure, Main/LCAL decrease, no lower-bound shortfall recovered & Bounded diagnostic & One point failure recovered; no lower-bound recovery or formal guarantee \\
\bottomrule
\end{tabular}

}

\end{table}%

The Main baseline selected method is the persistence-anchor Mondrian
conformal diagnostic. On the pooled-domain, cell-held-out task it
reaches coverage 94.8\% with CI 94.2\%-97.1\%, mean width 0.038, and
sample size 2323. The ablation package performs a separate
validation-only selection: its full-feature reference selects the
persistence-anchor standard conformal diagnostic, with coverage 92.6\%
and mean width 0.027. Removing the recent-SOH feature changes that
within-package selection to the gradient-boosting domain-Mondrian
conformal diagnostic, with coverage 80.8\% and mean width 0.221;
relative to the stated ablation reference, the coverage change is -11.8
pp and the width increase is 0.193. These deltas therefore are not
differences from the separately selected Main baseline. Figure
\ref{fig-evidence-flow} summarizes the resulting audit logic: raw
coverage is interpreted only after the measurement conditions are
examined.

\begin{figure}[!htbp]
\centering
\includegraphics[width=0.95\linewidth]{figures/evidence-flow.pdf}
\caption{The six-gate benchmark-integrity audit converts an unmet evidence requirement into a narrower admissible claim rather than a stopped analysis. Case verdicts are shown inside the gates; dashed orange or vermilion borders mark conditional, bounded-diagnostic, or failed gates. Pass* means that G5's reporting requirements are met while its serviceability interpretation remains bounded.}
\label{fig-evidence-flow}
\end{figure}

The cross-protocol diagnostics keep the result from becoming a simple
method-success story. Same-split QR/CQR/NGBoost comparisons show that
coverage can change only by changing the uncertainty family or accepting
wider intervals; for example, in the main held-out task, CQR is
coverage-aligned but has 2.72 times the persistence-anchor width. In the
common-cycle, no-recent-label audit, coverage is 92.9\% on Main, 29.5\%
on LNAS, and 3.1\% on LOxf. These values demonstrate heterogeneity, not
a universal monotone benefit of recent labels. On Main, the
clean-original no-recent schema attains a coverage of 92.4\%. Because
the as-used no-recent schema is flagged as contaminated in two tasks,
conclusions are based on the row-matched common-schema contrasts rather
than the contaminated hard-regime rows. Ancillary host-timing and
residual-bias screens are robustness checks only; they are not
interpreted as hardware or physical-degradation evidence.

The method-matched recent-label delta is gradient-boosting CQR coverage
under the common no-recent-SOH schema minus coverage with the recent-SOH
reference, evaluated on identical task rows. This comparison is
heterogeneous across tasks, which is central to the measurement-validity
interpretation. Table \ref{tbl-schema-deltas} reports it beside the
as-used schema-bundle contrast and its independent permuted-feature
control.

\begin{table}

\caption{\label{tbl-schema-deltas}Task-wise coverage sensitivity (pp; Bonferroni-adjusted cell-block normal intervals). Cells/rows gives effective cells and test records. Recent is no-recent minus recent-SOH; Bundle is as-used minus clean; Control is the permuted-feature contrast. G4 fails for any non-null control endpoint.}

\centering{

\centering
\scriptsize


\begin{tabular}{lllllcc}
\toprule
Task & Cells/rows & No-recent $-$ recent & Bundle & Control & Flag & G4 control \\
\midrule
Main & 10/2323 & +1.3 [-18.7, +21.4] & -4.1 [-21.6, +13.4] & -1.1 [-3.0, +0.7] & No & Pass \\
LCAL & 12/11592 & -10.1 [-16.8, -3.3] & +3.2 [+1.2, +5.1] & -0.6 [-1.0, -0.2] & Yes & Fail \\
TCAL & 12/5167 & -60.4 [-73.6, -47.2] & -2.2 [-6.9, +2.6] & -0.0 [-2.7, +2.6] & No & Fail \\
LNAS & 25/1728 & -31.9 [-40.7, -23.1] & -12.0 [-20.5, -3.4] & -1.0 [-2.0, +0.1] & Yes & Pass \\
TNAS & 25/783 & -41.5 [-73.4, -9.7] & -10.0 [-31.2, +11.2] & -0.1 [-1.6, +1.3] & No & Pass \\
LOxf & 8/511 & -37.6 [-53.1, -22.0] & +0.4 [-0.5, +1.3] & -0.2 [-1.0, +0.6] & No & Pass \\
TOxf & 8/233 & +4.3 [-1.6, +10.2] & +0.0 [+0.0, +0.0] & -0.4 [-2.1, +1.2] & No & Pass \\
\bottomrule
\end{tabular}

}

\end{table}%

The Control column in Table \ref{tbl-schema-deltas} displays the
coverage endpoint. TCAL's G4 control failure instead occurs at the
illustrative SOH cutoff 0.85, where the clean-original minus fixed-seed
permuted-control false-acceptance contrast is -26.4 pp {[}-52.1, -0.8{]}
after Bonferroni adjustment. The largest method-matched recent-label
contrast occurs in TCAL, where removing the recent-SOH reference changes
coverage by 60.4 pp; LNAS changes by 31.9 pp.~Across the six
cross-dataset tasks, five adjusted intervals support coverage
reductions; the TOxf estimate is positive but its interval includes
zero, as does Main's. The evidence therefore establishes task-dependent
magnitudes, not a supported positive reversal. Because these comparisons
hold the conformal family and task rows fixed, they isolate
information-set sensitivity more cleanly than the selection-changing
Main ablation described above. The Main ablation changes -11.8 pp, but
its selection rule switches conformal families and therefore measures
the response of the complete selection pipeline, not a same-method
feature effect.

The as-used schema-bundle contrast passes the experiment-level coverage
flag in 2 tasks: LCAL, LNAS. LNAS changes by -12.0 pp; its maximum
standardized mean difference (SMD) falls from 7.15 in the raw schema to
1.32 in the common schema. However, the as-used bundle jointly changes
\texttt{source\_row\_id}, \texttt{target\_cycle\_index}, and the
permuted control column, while the independent permuted-feature control
is also non-null in LCAL, TCAL. The documented claim gate therefore
returns false. The contrast cannot isolate a bookkeeping-field effect.
We interpret it only as a schema-bundle sensitivity warning that may
combine schema identity, finite-cell structure, and residual shift. This
negative-control failure blocks the stronger original artifact-mechanism
story.

LNAS shows a multiplicity-adjusted schema-bundle false-acceptance
increase of +100.0 pp at SOH threshold 0.90: the rate is 100.0\% under
the as-used schema and 0.0\% under the clean schema over 1127 actually
threshold-unacceptable records and 21 effective cells (exact paired-cell
p = 9.54e-07). The bundle's decision effect is directionally
heterogeneous. The largest LCAL contrast is -100.0 pp at SOH 0.65 (0.0\%
as-used versus 100.0\% clean; 1,936 records, 12 effective cells and 12
discordant cells; exact paired-cell p = 0.000488), which is in the
risk-decreasing direction and does not survive the corrected criterion.
Only LNAS passes correction, so the bundle cannot be described as
uniformly risk-increasing. Table \ref{tbl-false-acceptance-deltas}
reports the largest observed eligible contrast for each task, and the
complete 245-contrast ledger remains in the machine-readable results
archive. These are bundle-level decision associations, not isolated
bookkeeping effects or deployment-safety rates. Absolute decision risk
must therefore also be read directly from Table
\ref{tbl-decision-tradeoff}. Figure \ref{fig-coverage-deltas} visualizes
the task-dependent information-set and schema-bundle sensitivities.

\begin{table}

\caption{\label{tbl-false-acceptance-deltas}Largest eligible schema-bundle false-acceptance contrast by task (pp). As-used/clean gives the two base rates (percent); cells/rows gives the effective cell blocks and actually threshold-unacceptable records; discordant is the number of nonzero paired-cell differences. Intervals are Bonferroni-adjusted, and zero-SE cells show exact paired-cell p-values. Only LNAS passes correction.}

\centering{

\centering
\scriptsize


\begin{tabular}{llllllll}
\toprule
Task & Delta [CI] & As-used/clean & SOH & Cells/rows & Discordant & Exact p & Corrected? \\
\midrule
Main & +4.4 [-35.9, +44.6] & 6.0/1.7 & 0.75 & 7/712 & 4 & -- & No \\
LCAL & -100.0 & 0.0/100.0 & 0.65 & 12/1936 & 12 & 0.000488 & No \\
TCAL & +21.5 [-25.9, +68.8] & 28.6/7.1 & 0.85 & 12/5036 & 9 & -- & No \\
LNAS & +100.0 & 100.0/0.0 & 0.90 & 21/1127 & 21 & 9.54e-07 & Yes \\
TNAS & +45.8 [-19.8, +100.0] & 97.6/51.7 & 0.75 & 14/288 & 8 & -- & No \\
LOxf & +0.0 & 0.0/0.0 & 0.90 & 8/342 & 0 & 1 & No \\
TOxf & +0.0 & 0.0/0.0 & 0.90 & 8/233 & 0 & 1 & No \\
\bottomrule
\end{tabular}

}

\end{table}%

\begin{figure}[!htbp]
\centering
\includegraphics[width=0.9\linewidth]{figures/coverage-deltas.pdf}
\caption{Per-task recent-label and as-used schema-bundle coverage deltas (percentage points) with Bonferroni-corrected normal intervals. Recent-label deltas are large and negative in several tasks; bundle deltas are smaller, with LCAL and LNAS meeting the coverage flag. The endpoint-wide negative-control gate fails in LCAL and TCAL. The dashed line marks zero.}
\label{fig-coverage-deltas}
\end{figure}

Decision-utility rates are reported with denominators and cell-block
intervals rather than as unqualified safety claims. Here, a
false-acceptance event occurs when the actual SOH falls below the
serviceability threshold but the predicted interval incorrectly implies
threshold acceptability, while a false-rejection event occurs when an
actually threshold-acceptable SOH is incorrectly flagged as
threshold-unacceptable. Table \ref{tbl-decision-tradeoff} reports
representative SOH threshold decision cells at threshold 0.80. The table
illustrates why coverage alone is insufficient for a threshold-decision
interpretation: the leave-NASA CQR cell has a lower false-acceptance
point estimate but leaves a large uncertain region, whereas NGBoost has
a high false-acceptance diagnostic. These thresholds and cells are
illustrative boundary evidence, not deployment-safety rates; application
use requires threshold-specific validation.

\begin{table}

\caption{\label{tbl-decision-tradeoff}Illustrative threshold-0.80 decision cells (percent; uncorrected 95\% cell-block CI). Here c/n gives effective cells/records; BD marks a bootstrap-degenerate all-zero boundary, and Sparse denotes any denominator below 5.}

\centering{

\centering
\scriptsize


\begin{tabular}{ll>{\raggedright\arraybackslash}p{0.21\linewidth}>{\raggedright\arraybackslash}p{0.21\linewidth}>{\raggedright\arraybackslash}p{0.21\linewidth}c}
\toprule
Task & Method & False acceptance & False rejection & Uncertain & Flag \\
\midrule
Main & Persistence anchor & 1.4 [0.1, 2.7]; c=9, n=1270 & 2.1 [1.1, 3.1]; c=10, n=1053 & 18.1 [12.0, 24.1]; c=10, n=2323 & OK \\
Main & CQR & 1.1 [0.0, 2.4]; c=9, n=1270 & 0.7 [0.0, 1.7]; c=10, n=1053 & 30.9 [12.4, 49.4]; c=10, n=2323 & OK \\
Main & NGBoost & 1.0 [0.0, 2.3]; c=9, n=1270 & 6.9 [0.0, 19.2]; c=10, n=1053 & 33.3 [25.3, 41.3]; c=10, n=2323 & OK \\
LNAS & Persistence anchor & 1.9 [0.3, 3.5]; c=16, n=627 & 1.0 [0.1, 1.9]; c=25, n=1101 & 6.1 [3.9, 8.4]; c=25, n=1728 & OK \\
LNAS & CQR & 1.4 [0.1, 2.7]; c=16, n=627 & 0.0 [BD]; c=25, n=1101 & 39.6 [28.2, 51.0]; c=25, n=1728 & OK \\
LNAS & NGBoost & 69.2 [46.8, 91.6]; c=16, n=627 & 0.0 [BD]; c=25, n=1101 & 11.5 [1.1, 21.8]; c=25, n=1728 & OK \\
\bottomrule
\end{tabular}

}

\end{table}%

All representative G5 cells in Table \ref{tbl-decision-tradeoff} report
both error directions and exceed the denominator-5 suppression floor.
The two zero-event LNAS false-rejection cells are marked
bootstrap-degenerate because cell resampling cannot estimate
unseen-event uncertainty from an all-zero boundary. The gate therefore
passes its reporting requirement, while the application meaning remains
bounded because the thresholds are illustrative, the effective cell
counts are limited, and zero observed events do not establish zero
operational risk.

The localized comparator is below the nominal point target in Main,
LCAL, TCAL, and LOxf. In LCAL, the standard comparator has coverage
75.5\%, and localization changes coverage by -8.3 pp without repairing
that failure. In Main, the standard comparator is above target at 95.9\%
with lower bound 95.0\%; localization changes coverage by -6.5 pp and
lowers the bound to 83.7\%, thereby converting a pass into a
point-coverage and lower-bound shortfall. TNAS is the only task in which
localization converts a standard point-coverage failure into a pass;
LNAS and TOxf already meet the point target under the standard
comparator. No lower-bound shortfall crosses the nominal boundary from
failure to pass. Because the comparator uses distance-based residual
localization rather than density-ratio weighting, these are empirical
diagnostic outcomes without a covariate-shift coverage guarantee. Table
\ref{tbl-adaptive-envelope} is therefore a sensitivity envelope, not
evidence of conformal superiority.

\begin{table}

\caption{\label{tbl-adaptive-envelope}Heuristic localized-residual CQR coverage by task. Cells/rows gives the effective battery-cell blocks and test records. Std (low) and Local (low) give point coverage with the 95\% cell-block lower bound in parentheses; Delta is local minus standard coverage (pp). Pt $\geq$ nom marks nominal point coverage, not a formal coverage guarantee.}

\centering{

\centering
\footnotesize


\begin{tabular}{llllcc}
\toprule
Task & Cells/rows & Std (low) & Local (low) & Delta & Pt $\geq$ nom \\
\midrule
Main & 10/2323 & 95.9\% (95.0\%) & 89.4\% (83.7\%) & -6.5 pp & N \\
LCAL & 12/11592 & 75.5\% (69.1\%) & 67.2\% (60.9\%) & -8.3 pp & N \\
TCAL & 12/5167 & 65.6\% (52.9\%) & 68.0\% (56.5\%) & +2.4 pp & N \\
LNAS & 25/1728 & 90.4\% (86.0\%) & 91.2\% (87.0\%) & +0.8 pp & Y \\
TNAS & 25/783 & 88.0\% (83.6\%) & 93.7\% (89.7\%) & +5.7 pp & Y \\
LOxf & 8/511 & 42.9\% (35.0\%) & 42.9\% (35.0\%) & +0.0 pp & N \\
TOxf & 8/233 & 98.7\% (96.9\%) & 98.7\% (96.9\%) & +0.0 pp & Y \\
\bottomrule
\end{tabular}

}

\end{table}%

\section{Discussion}\label{sec-discussion}

The audit changes the interpretation of cross-dataset conformal SOH
benchmarks. Coverage is conditional on the benchmark's measurement
validity and information set; it is not automatically a reliability or
serviceability result. Recent-label access can make a benchmark easier
than the intended forecasting problem. The as-used schema-bundle
contrast reveals an unstable evaluation design, but it jointly changes
multiple columns, and the endpoint-wide independent-control gate fails
in LCAL and TCAL. The corrected LNAS false-acceptance association is
therefore retained at the bundle level without attribution to
bookkeeping fields or promotion to a causal schema mechanism. Residual
drift, finite-cell variation, and label-definition differences remain
plausible contributors.

The protocol's position relative to the current state of the art is
complementary rather than competitive. Existing conformal work in
prognostics and battery state estimation --- split-conformal RUL
pipelines, robust online conformal prediction, conformalized quantile
regression for electrified transportation --- concentrates on producing
better intervals
\citep{xu2026uncertaintyaware, wang2025robust, javanmardi2023conformal, piao2025crulp, soon2026enhancing};
the transfer and domain-generalization literature concentrates on making
models survive shift
\citep{nejjar2026uncertaintyguided, fei2026decoupled}. Neither line
audits whether the benchmark that certifies those intervals measures the
intended deployment question. The coverage collapses observed here are
the empirical face of known theory --- coverage degradation under
non-exchangeability is bounded, not prevented
\citep{barber2023beyond, oliveira2024split}, and exchangeability
violations are already documented in industrial time series
\citep{zhang2025uncertainty} --- but the theory does not tell an analyst
which of these failure channels a specific benchmark contains. The
protocol does: it separates information-set effects from schema-bundle
sensitivity and residual drift with governed component gates. Adaptive
and online conformal methods do not directly rescue this setting either,
because a one-shot cross-dataset evaluation provides no sequential
feedback for their updates
\citep{gibbs2021adaptive, gibbs2024conformal}. The one-task
point-coverage recovery achieved by the localized-residual comparator is
consistent with the localized-conformal literature, in which reweighted
calibration improves conditional behavior only where local residual
structure is informative and offers no general conditional guarantee
\citep{guan2022localized, hore2024conformal}.

For reliability and battery-health practitioners, the engineering
implication is a pre-submission benchmark-integrity discipline.
Prognostic intervals are consumed by predictive-maintenance planning
\citep{mitici2023dynamic, pedersen2022optimizing}, condition-based
policies \citep{zhuang2023prognostic, depater2021predictive}, and
multi-component scheduling
\citep{wang2026dynamic, wang2025multiobjective}; battery-health
predictions are also used to optimize maintenance actions
\citep{wang2026multiaction}. The observed +100.0 pp false-acceptance
change at an illustrative serviceability threshold is therefore not
merely a coverage statistic: if such a benchmark were used unaudited to
set a maintenance threshold, the resulting decision evidence would be
mispriced. Quantifying that mispricing for a fleet requires
application-specific threshold validation, which this study does not
perform. Moreover, the opposite-direction -100.0 pp LCAL contrast shows
that the bundle is not uniformly risk-increasing. Analysts should
declare whether a recent measured SOH is available at prediction time,
demonstrate task-appropriate split integrity, use cell-aware
uncertainty, test schema claims against independent negative controls,
and report threshold decisions with effective cell counts and record
denominators. A localization comparator can test sensitivity, but it
cannot substitute for a valid conformal guarantee. The audit protocol of
Table \ref{tbl-audit-protocol} condenses this discipline into six
governed gates; each has an evidentiary requirement and a failure
consequence, so a failure bounds the claim rather than terminating the
analysis. In this respect the protocol extends the
reliability-engineering practice of assessing machine-learning evidence
before use
\citep{zio2022prognostics, chen2024evaluating, panic2027simple, salinascamus2025comprehensive}
to conformal SOH benchmarks.

The recent-label finding is deliberately not stated as a monotone
benefit. Method-matched deltas are strongly negative in TCAL, LNAS,
TNAS, and LOxf, smaller but still negative in LCAL, near zero in Main,
and positive in TOxf. That reversal matters: recent-SOH is not simply a
universally helpful covariate; it changes the estimand and interacts
with the target task. The practical failure is therefore undeclared
information availability --- a failure mode the label-scarce deployment
literature makes concrete, since in-service capacity labels are exactly
what adaptive battery prognostics cannot assume
\citep{zhang2022adaptive, pan2024probabilistic}, and
feature-availability choices are already known to mislead
battery-lifetime prediction \citep{geslin2023selecting}. A benchmark
using an oracle recent label cannot support a no-recent-label deployment
claim without a separate evaluation.

The localized-residual comparator sharpens the interpretation without
turning the study into a method-success paper. It changes empirical
coverage in several tasks and recovers exactly one standard
point-coverage failure, but other tasks remain below target. Because it
is not a density-ratio conformal method, its results show only that
calibration is sensitive to local residual composition; they do not
establish shift-robust coverage.

The main limitation is external breadth. The evidence base is limited to
CALCE, NASA, and Oxford, so the demonstrated claim is not that every
battery SOH benchmark has the same sensitivity. A Massachusetts
Institute of Technology (MIT)-Stanford-Toyota fast-charge dataset
\citep{severson2019datadriven} is excluded because its
lithium-iron-phosphate chemistry and fast-charge-to-failure protocol
would add chemistry and protocol confounding, and its records did not
pass the same preprocessing and schema-compatibility checks. Broader
modern and field datasets are needed for external validation.

Several design limitations follow directly from the audit framing.
Calibration residuals are cycle-level observations from longitudinal
cell histories; cell-block intervals quantify rate uncertainty but do
not create independent exchangeable histories. A stronger follow-up
should use cell-disjoint prospective splits, report row-weighted and
cell-macro coverage, and show worst-cell behavior near serviceability
thresholds. The article treats recent-SOH as an oracle recent label;
deployment-facing work should compare oracle, estimated, and unavailable
recent-SOH on identical rows. The capacity-based label is harmonized
only through the current preprocessing interface. G4's independent
control uses a single permutation draw (\(B=1\)), so it cannot
characterize a permutation null distribution. Because the control acts
only as a veto, a non-null draw can withhold mechanism attribution but
cannot create it; repeated permutation draws would quantify how often a
null feature produces a control effect as large as the LCAL result.
Likewise, the zero-event false-rejection cells in Table
\ref{tbl-decision-tradeoff} are bootstrap-degenerate boundaries rather
than evidence of zero underlying risk. Moreover, the original frozen
processed CSV could not be recovered byte-for-byte; the repaired
analyses were regenerated from the preserved raw holdings under Python
3.11.15 and scikit-learn 1.7.2. Although raw split membership was
reproduced, this provenance discontinuity limits exact numerical
replication of the earlier exploratory run and is one reason the present
conclusions rely only on the corrected rerun.

The serviceability-threshold language is intentional. The audited SOH
thresholds (0.65, 0.70, 0.75, 0.80, 0.85, 0.90) are illustrative
capacity-SOH cutoffs used to demonstrate the diagnostic quantities.
Their use in maintenance, warranty, or end-of-life decisions requires
application-specific validation; this manuscript does not establish them
as electrochemical safety boundaries or thermal-runaway risk thresholds.
Additional field data, prospective forecasting splits, and
application-specific risk models are needed before the audit can be
treated as a general battery-management validation protocol.

\section{Conclusion}\label{sec-conclusion}

This paper develops and demonstrates a governed benchmark-integrity
audit protocol for conformal battery SOH evidence. In the corrected
real-data analyses, method-matched recent-label removal produces
statistically supported coverage reductions in five cross-dataset tasks,
with task-dependent magnitudes, while Main and TOxf remain compatible
with no effect. The protocol also detects that the as-used schema bundle
changes coverage in two tasks and increases LNAS false acceptance at one
corrected threshold. The LCAL decision contrast has the opposite sign
and does not pass multiplicity correction. Because the bundle jointly
changes multiple columns and the endpoint-wide independent control is
non-null, the protocol reports these results at the bundle level only
--- demonstrating how a governed control prevents schema sensitivity
from being over-interpreted as a mechanism.

The five contributions remain aligned with that evidence. First, the
paper defines an explicit six-gate benchmark-integrity protocol. Second,
it quantifies task-dependent recent-label sensitivity with
method-matched contrasts. Third, it combines multiplicity-adjusted
schema contrasts with an independent negative control. Fourth, it
reports a corrected LNAS link between schema-bundle sensitivity and
threshold false acceptance. Fifth, it bounds the final diagnostic as a
localized-residual sensitivity envelope without a covariate-shift
guarantee.

Every gate specifies a condition, an evidentiary requirement, and a
bound on the admissible claim when the requirement is not met. The gates
are formulated independently of the three datasets, but their detection
performance and portability have not been validated; this case
demonstrates executability and claim bounding only. Because the six-gate
synthesis was abstracted from the diagnostic program it audits, it is
not an independent test of detection power. Applying the protocol
prospectively to benchmarks that the authors did not construct is the
necessary next validation step.

Within the three public laboratory datasets studied here, the
implication for reliability engineering is direct: coverage should
support a serviceability interpretation only after the benchmark's
information set, split integrity, uncertainty units, controls, and
decision consequences have been audited. Broader field data, prospective
forecasting splits, and repeated-control designs remain necessary before
generalizing beyond this case.

\section*{Data Availability}\label{data-availability}
\addcontentsline{toc}{section}{Data Availability}

This study uses real cycle-level SOH records from three public battery
datasets: the CALCE Battery Research Group battery data repository
\citep{calce_battery_data}, the NASA Ames Prognostics Data Repository
\citep{nasa_pcoe_battery_data}, and the Oxford Battery Degradation
Dataset \citep{howey2017oxford}, all accessed 2026-06-02. The versioned
derived-data package, source-data provenance, split manifests, run
manifests, archived analysis plans, and regeneration code underlying
every reported value are publicly available at the immutable replication
commit
https://github.com/QuantDiviner/NEX-CP-conformal-soh-audit/tree/42f1120ed1f5cbc0282a706ebd4bec24a8d19d6e.
Supplementary Tables S1--S2 provide the case-implementation and
provenance matrix.

\section*{Code Availability}\label{code-availability}
\addcontentsline{toc}{section}{Code Availability}

The complete analysis code, locked dependencies, and run manifests are
available at the immutable replication commit above. The repository
regenerates the scripted results table from which all manuscript values
are drawn.

\section*{Declaration of Competing
Interest}\label{declaration-of-competing-interest}
\addcontentsline{toc}{section}{Declaration of Competing Interest}

The authors declare that they have no known competing financial
interests or personal relationships that could have appeared to
influence the work reported in this paper.

\section*{Funding}\label{funding}
\addcontentsline{toc}{section}{Funding}

This work was supported by the National Social Science Foundation of
China (grant no. 25BTJ045) and the Natural Science Foundation of Jiangxi
Province (grant no. 20232BAB201025), both awarded to Q.L. The funders
had no role in study design, data collection and analysis, decision to
publish, or preparation of the manuscript.

\section*{Declaration of Generative AI and AI-Assisted Technologies in
the Manuscript Preparation
Process}\label{declaration-of-generative-ai-and-ai-assisted-technologies-in-the-manuscript-preparation-process}
\addcontentsline{toc}{section}{Declaration of Generative AI and
AI-Assisted Technologies in the Manuscript Preparation Process}

During the preparation of this work, the authors used OpenAI Codex for
code editing assistance, manuscript formatting, and language polishing.
The authors subsequently reviewed and edited the content, verified all
reported evidence and results, and take full responsibility for the
content of the publication.

\section*{CRediT Authorship Contribution
Statement}\label{credit-authorship-contribution-statement}
\addcontentsline{toc}{section}{CRediT Authorship Contribution Statement}

\textbf{Qingsong Shan}: Software, Data curation, Formal analysis,
Validation, Writing - Review \& Editing. \textbf{Qianning Liu}:
Conceptualization, Methodology, Investigation, Visualization, Writing -
Original Draft, Writing - Review \& Editing, Supervision, Funding
acquisition.

\section*{Corresponding Author}\label{corresponding-author}
\addcontentsline{toc}{section}{Corresponding Author}

Qianning Liu\\
School of Statistics and Data Science\\
Key Laboratory of Data Science in Finance and Economics\\
Jiangxi University of Finance and Economics\\
Nanchang, China\\
Email: liuqianning@jxufe.edu.cn


\bibliography{references/main.bib}



\end{document}
